%% ============================================================
%%  Aula 7 — Gradientes de política II
%%  Adaptação em pt-BR do CS234 (Stanford, Emma Brunskill), Lecture 7
%%  Prof. Marcelo Keese Albertini — Universidade Federal de Uberlândia
%%  Compilar com XeLaTeX (2x) — latexmk -xelatex aula07.tex
%% ============================================================
\documentclass[aspectratio=169, 10pt]{beamer}
%% .sty do projeto na raiz do repositório (../)
\makeatletter
\def\input@path{{./}{../}}
\makeatother


\usetheme{AprendizadoRL}      % ANTES do babel — fontspec precisa vir primeiro
\usepackage{babel}
\babelprovide[import, main]{brazilian}
\usepackage{booktabs}
\usepackage{cancel}
\usepackage{etoolbox}

%% ---- caixas com geometria compacta (mesma identidade, menos respiro) ----
\tcbset{rlcaixaC/.style={
  enhanced, boxrule=0pt, frame hidden, arc=1.6mm,
  left=2mm, right=2mm, top=0.7mm, bottom=0.8mm,
  boxsep=0.6mm, before skip=0.7mm, after skip=0.8mm,
  borderline west={2.4pt}{0pt}{#1},
  colback=rlFundo, fonttitle=\bfseries\small, coltitle=#1,
  attach title to upper={\par\vspace{0.2ex}}, title after break={}}}
\renewtcolorbox{defbox}[2][]{rlcaixaC=rlTeal, title={Definição\ifblank{#2}{}{ — #2}}, #1}
\renewtcolorbox{teobox}[2][]{rlcaixaC=rlAzulClaro, title={#2}, colback=rlAzulClaro!6, #1}
\renewtcolorbox{algbox}[2][]{rlcaixaC=rlAzul, title={Algoritmo — #2}, colback=rlAzul!5, #1}
\renewtcolorbox{ideiabox}[1][]{rlcaixaC=rlAmbar, fontupper=\small,
  title={Intuição}, colback=rlAmbar!10, coltitle=rlLaranja, #1}
\renewtcolorbox{alertabox}[1][]{rlcaixaC=rlVermelho, fontupper=\small,
  title={Atenção}, colback=rlVermelho!5, #1}
\renewtcolorbox{respbox}[1][]{rlcaixaC=rlVerde, fontupper=\small,
  title={Resposta}, colback=rlVerde!7, #1}

%% ---- compactação global de listas (densidade de slide) ----
\AtBeginEnvironment{itemize}{\setlength{\itemsep}{0.18ex}\setlength{\parsep}{0pt}%
  \setlength{\topsep}{0.2ex}\setlength{\partopsep}{0pt}}
\AtBeginEnvironment{enumerate}{\setlength{\itemsep}{0.18ex}\setlength{\parsep}{0pt}%
  \setlength{\topsep}{0.2ex}\setlength{\partopsep}{0pt}}
\setlength{\abovedisplayskip}{4pt}
\setlength{\belowdisplayskip}{4pt}
\setlength{\abovedisplayshortskip}{2pt}
\setlength{\belowdisplayshortskip}{2pt}

\title[Aula 7 — PPO e GAE]{Aula 7: Gradientes de política II}
\subtitle[Aprendizagem por Reforço]{Amostragem por importância \textbullet\ PPO \textbullet\ GAE \textbullet\ Melhoria monotônica}
\author[Prof. Marcelo K. Albertini]{Prof. Marcelo Keese Albertini \\
  \footnotesize Universidade Federal de Uberlândia \\
  \footnotesize adaptado de Emma Brunskill, CS234 (Stanford)}
\institute{Universidade Federal de Uberlândia}
\date{2026}

% ---- Nota discreta: perguntas ficam para o professor conduzir em aula
\newcommand{\paradiscussao}{\smallskip
  {\footnotesize\color{rlCinza}\itshape Pergunta para discussão conduzida em aula.}}

%% ------------------------------------------------------------
%% Macros locais desta aula
%% ------------------------------------------------------------
\newcommand{\thv}{\theta}
\newcommand{\thk}{\thv_k}
\newcommand{\polth}{\pol_{\thv}}
\newcommand{\polk}{\pol_{k}}
\newcommand{\polnovo}{\pol'}
\newcommand{\gth}{\nabla_{\thv}}
\newcommand{\traj}{\tau}
\newcommand{\Vhat}{\hat{V}}
\newcommand{\Qhat}{\hat{Q}}
\newcommand{\Ahat}{\hat{A}}
\newcommand{\Adv}{A}
\newcommand{\Advpi}{A^{\pol}}
\newcommand{\Dkl}{D_{\mathrm{KL}}}
\newcommand{\dpi}{d^{\pol}}
\newcommand{\Lsub}[1]{\mathcal{L}_{#1}}
\DeclareMathOperator{\clip}{clip}
\DeclareMathOperator{\ent}{H}
\newcommand{\razao}{r_t(\thv)}
\newcommand{\Lclip}{\mathcal{L}^{\mathrm{CLIP}}}
\newcommand{\feat}{x}                    % vetor de atributos do estado
\newcommand{\mufeat}{\mu}                % expectativas de atributos
\newcommand{\wv}{w}
\newcommand{\Var}{\operatorname{Var}}
\newcommand{\Dtot}{\mathcal{D}}
\newcommand{\polE}{\pol^{*}}             % política do especialista

\begin{document}

\begin{frame}[plain, noframenumbering]\titlepage\end{frame}

%% ------------------------------------------------------------
\begin{frame}{Onde estamos}{Da estimativa do gradiente ao tamanho do passo}

  \begin{itemize}\setlength{\itemsep}{0.45ex}
    \item \textbf{Aulas 2--3.} Com um modelo do mundo, calcular $\Vstar$ e $\polstar$
      por programação dinâmica.
    \item \textbf{Aula 4.} Sem modelo, aprender $\Qhat(s,a;w)$ a partir de amostras.
    \item \textbf{Aulas 5--6.} Parametrizar a política e subir o gradiente de
      $J(\polth)$; medir distância entre políticas no espaço de \emph{distribuições}.
    \item \textbf{Hoje.} Fechar a família: como reaproveitar dados de uma política
      velha (\alert{amostragem por importância}), como o \alert{PPO} torna isso
      barato, como o \alert{GAE} equilibra viés e variância na vantagem, e
      \emph{por que} tudo isso funciona (teoria da melhoria monotônica).
  \end{itemize}

  \begin{ideiabox}
    A outra metade --- aprender com \alert{demonstrações}, sem recompensa escrita:
    clonagem de comportamento, DAgger e RL inverso --- é o assunto da Aula 8.
  \end{ideiabox}
\end{frame}

%% ------------------------------------------------------------
\begin{frame}{Plano de hoje}

  \begin{enumerate}\setlength{\itemsep}{0.5ex}
    \item \textbf{Amostragem por importância} --- a ferramenta que permite avaliar
      uma política nova com dados de outra, e por que ela sozinha não basta.
    \item \textbf{Otimização proximal de política (PPO)} --- penalidade KL adaptativa
      e objetivo recortado.
    \item \textbf{Estimador generalizado de vantagem (GAE)} --- um botão único,
      $\lambda$, que percorre todo o espectro entre TD(0) e Monte Carlo.
    \item \textbf{Teoria da melhoria monotônica} --- o limite de desempenho relativo
      e o argumento de \emph{majorizar--maximizar} que justifica TRPO e PPO.
  \end{enumerate}
\end{frame}

%% ============================================================
\section{Reaproveitar dados: amostragem por importância}
%% ============================================================

%% ------------------------------------------------------------
\begin{frame}{O que ainda incomoda no gradiente de política}

  Os algoritmos de gradiente de política resolvem
  \[
    \max_{\thv} \; J(\polth) \defeq \E_{\traj \sim \polth}
      \Bigl[\textstyle\sum_{t=0}^{\infty} \gamma^t r_t\Bigr],
    \qquad
    g = \gth J(\polth) = \E_{\traj\sim\polth}
      \Bigl[\textstyle\sum_{t\ge 0}\gamma^t \gth \log \polth(a_t\mid s_t)\,\Advpi(s_t,a_t)\Bigr].
  \]

  \begin{itemize}\setlength{\itemsep}{0.35ex}
    \item \textbf{Eficiência amostral ruim.} O estimador é \emph{on-policy}: cada
      lote de trajetórias serve a \emph{uma} atualização e é descartado.
    \item \textbf{Distância em $\thv$ $\neq$ distância em $\pol$.} O passo é dado
      no espaço de parâmetros, mas o que estraga (ou salva) o aprendizado é o
      deslocamento no espaço de \emph{políticas}.
  \end{itemize}

  \begin{defbox}{Espaço de políticas}
    No caso tabular, o conjunto de matrizes estocásticas
    \[
      \Pi = \Bigl\{ p \in \R^{\abs{\gS}\times\abs{\gA}} \;:\;
        \textstyle\sum_a p_{sa} = 1,\; p_{sa}\ge 0 \Bigr\}.
    \]
    A parametrização $\thv \mapsto \polth$ é um mergulho de $\R^n$ em $\Pi$ --- e
    esse mergulho \emph{distorce distâncias}.
  \end{defbox}
\end{frame}

%% ------------------------------------------------------------
\begin{frame}{O mesmo passo em $\thv$, dois deslocamentos muito diferentes}

  Política gaussiana $\polth(a\mid s) = \mathcal{N}(\mu_\thv(s), \sigma^2)$.
  Desloque $\mu$ da mesma quantidade $\Delta\mu = 0{,}5$ nos dois casos:

  \begin{center}
  \scalebox{0.92}{\begin{tikzpicture}[font=\scriptsize]
    % ---- painel esquerdo: sigma pequeno
    \begin{scope}
      \draw[->, rlCinza] (-0.3,0) -- (4.4,0) node[right] {$a$};
      \draw[->, rlCinza] (-0.3,0) -- (-0.3,2.0);
      \draw[rlAzul, very thick, domain=0:4, samples=90]
        plot (\x, {1.75*exp(-(\x-1.4)^2/(2*0.045))});
      \draw[rlLaranja, very thick, domain=0:4, samples=90]
        plot (\x, {1.75*exp(-(\x-1.9)^2/(2*0.045))});
      \node[align=center, rlAzul] at (2.05,-0.45) {$\sigma = 0{,}21$};
      \node[align=center, text=rlVermelho, font=\scriptsize\bfseries] at (2.05,-0.92)
        {políticas quase disjuntas};
      \draw[<->, rlCinza, thick] (1.4,1.95) -- (1.9,1.95) node[midway, above] {$\Delta\mu$};
    \end{scope}
    % ---- painel direito: sigma grande
    \begin{scope}[xshift=6cm]
      \draw[->, rlCinza] (-0.3,0) -- (4.4,0) node[right] {$a$};
      \draw[->, rlCinza] (-0.3,0) -- (-0.3,2.0);
      \draw[rlAzul, very thick, domain=-0.2:4, samples=90]
        plot (\x, {1.75*exp(-(\x-1.4)^2/(2*0.9))});
      \draw[rlLaranja, very thick, domain=-0.2:4, samples=90]
        plot (\x, {1.75*exp(-(\x-1.9)^2/(2*0.9))});
      \node[align=center, rlAzul] at (2.05,-0.45) {$\sigma = 0{,}95$};
      \node[align=center, text=rlVerde!70!black, font=\scriptsize\bfseries] at (2.05,-0.92)
        {políticas quase idênticas};
      \draw[<->, rlCinza, thick] (1.4,1.95) -- (1.9,1.95) node[midway, above] {$\Delta\mu$};
    \end{scope}
  \end{tikzpicture}}
  \end{center}

  \vspace{0.2em}
  \begin{alertabox}
    Um único passo $\alpha$ não pode estar certo nos dois regimes. É por isso que
    a taxa de aprendizado do gradiente de política é notoriamente difícil de
    ajustar --- e por que faz sentido controlar o passo pela \alert{divergência KL}
    entre as políticas, não pela norma em $\thv$.
  \end{alertabox}
\end{frame}

%% ------------------------------------------------------------
\begin{frame}{Amostragem por importância}

  \begin{defbox}{Identidade de mudança de medida}
    Para estimar uma esperança sob $P$ usando amostras de $Q$ (com $Q(x)>0$ sempre
    que $P(x)f(x)\neq 0$):
    \[
      \E_{x\sim P}[f(x)] \;=\; \E_{x\sim Q}\!\left[\frac{P(x)}{Q(x)} f(x)\right]
      \;\approx\; \frac{1}{\abs{\Dtot}}\sum_{x\in\Dtot}\frac{P(x)}{Q(x)} f(x),
      \qquad \Dtot \sim Q .
    \]
    A razão $P(x)/Q(x)$ é o \alert{peso de importância} de $x$.
  \end{defbox}

  \begin{itemize}\setlength{\itemsep}{0.3ex}
    \item A identidade é exata --- o estimador é \textbf{não-viesado} para qualquer $Q$
      admissível.
    \item O problema nunca é o viés. É a \textbf{variância}.
  \end{itemize}

  \begin{ideiabox}
    Intuição: reponderamos cada amostra pelo quanto ela era \emph{mais provável sob
    $P$ do que sob $Q$}. Se $Q$ raramente visita uma região onde $P f$ é grande,
    as poucas amostras que caem lá carregam um peso enorme --- e o estimador passa
    a depender delas.
  \end{ideiabox}
\end{frame}

%% ------------------------------------------------------------
\begin{frame}{A variância do estimador de importância}

  Com $N$ amostras i.i.d.\ de $Q$, escrevendo $\hat\mu_Q$ para o estimador:
  \begin{align*}
    \Var(\hat\mu_Q)
      &= \frac{1}{N}\Var_{x\sim Q}\!\left[\frac{P(x)}{Q(x)}f(x)\right]
       = \frac{1}{N}\left(
          \E_{x\sim Q}\!\left[\frac{P(x)^2}{Q(x)^2}f(x)^2\right]
          - \E_{x\sim Q}\!\left[\frac{P(x)}{Q(x)}f(x)\right]^{2}\right)\\[0.3em]
      &= \frac{1}{N}\left(
          \mdestaque{rlVermelho}{\E_{x\sim P}\!\left[\frac{P(x)}{Q(x)}f(x)^2\right]}
          - \E_{x\sim P}[f(x)]^{2}\right).
  \end{align*}

  \begin{alertabox}
    O termo destacado é o vilão. Ele pesa $f(x)^2$ pela razão $P/Q$: se $P/Q$ for
    grande exatamente onde $f$ é grande, a variância explode --- e pode ser
    \emph{infinita} mesmo com $P$ e $Q$ perfeitamente bem-comportadas.
  \end{alertabox}

  \begin{itemize}
    \item Diagnóstico prático: o \emph{tamanho amostral efetivo}
      $\mathrm{ESS} = (\sum_i \omega_i)^2 / \sum_i \omega_i^2$, com $\omega_i = P(x_i)/Q(x_i)$.
      Se $\mathrm{ESS} \ll N$, poucas amostras estão fazendo todo o trabalho.
  \end{itemize}
\end{frame}

%% ------------------------------------------------------------
\begin{frame}{Amostragem por importância no gradiente de política}
  \framesubtitle{Comprimindo $\polth$ para $\thv$ na notação}

  \begin{align*}
    g = \gth J(\thv)
      &= \E_{\traj\sim\thv}\Bigl[\textstyle\sum_{t\ge0}\gamma^t \gth\log\polth(a_t\mid s_t) \Adv^{\thv}(s_t,a_t)\Bigr]\\
      &= \sum_{\traj}\sum_{t\ge0}\gamma^t P(\traj_t\mid \thv)\,\gth\log\polth(a_t\mid s_t)\Adv^{\thv}(s_t,a_t)\\
      &= \E_{\traj\sim\thv'}\Bigl[\textstyle\sum_{t\ge0}
          \mdestaque{rlLaranja}{\frac{P(\traj_t\mid\thv)}{P(\traj_t\mid\thv')}}
          \gamma^t\gth\log\polth(a_t\mid s_t)\Adv^{\thv}(s_t,a_t)\Bigr].
  \end{align*}

  \pause
  A razão colapsa: dinâmica e distribuição inicial aparecem nos dois lados.
  \[
    \frac{P(\traj_t\mid\thv)}{P(\traj_t\mid\thv')}
      = \frac{\cancel{\mu(s_0)}\prod_{t'=0}^{t}\cancel{P(s_{t'+1}\mid s_{t'},a_{t'})}\,\polth(a_{t'}\mid s_{t'})}
             {\cancel{\mu(s_0)}\prod_{t'=0}^{t}\cancel{P(s_{t'+1}\mid s_{t'},a_{t'})}\,\pol_{\thv'}(a_{t'}\mid s_{t'})}
      = \prod_{t'=0}^{t}\frac{\polth(a_{t'}\mid s_{t'})}{\pol_{\thv'}(a_{t'}\mid s_{t'})}.
  \]

  \begin{ideiabox}
    É o mesmo cancelamento da Aula 5: modelo do mundo desconhecido, mas irrelevante.
  \end{ideiabox}
\end{frame}

%% ------------------------------------------------------------
\begin{frame}{O produto é uma bomba-relógio}

  \[
    \frac{P(\traj_t\mid\thv)}{P(\traj_t\mid\thv')}
      = \prod_{t'=0}^{t}\frac{\polth(a_{t'}\mid s_{t'})}{\pol_{\thv'}(a_{t'}\mid s_{t'})}
  \]

  \begin{columns}[T]
  \begin{column}{0.52\textwidth}
    Suponha que em cada passo a razão valha $1{,}02$ --- políticas praticamente
    indistinguíveis, $2\%$ de diferença por ação.

    \medskip
    \begin{center}\small
    \begin{tabular}{@{}rr@{}}
      \toprule
      $t$ & $\prod$ razões \\
      \midrule
      \phantom{0}10 & $1{,}22$ \\
      \phantom{0}50 & $2{,}69$ \\
      100 & $7{,}24$ \\
      500 & $2{,}0\times 10^{4}$ \\
      \bottomrule
    \end{tabular}
    \end{center}
  \end{column}
  \begin{column}{0.46\textwidth}
    \begin{alertabox}
      Diferenças minúsculas \emph{multiplicadas} ao longo do horizonte viram uma
      diferença enorme. Na direção oposta, o peso vai a zero e a trajetória some
      do estimador.
    \end{alertabox}

    \begin{ideiabox}
      \textbf{A pergunta central desta aula:}
      como aproveitar de forma eficiente os dados que já temos da política velha,
      sem pagar o preço da razão de importância sobre trajetórias inteiras?
    \end{ideiabox}
  \end{column}
  \end{columns}
\end{frame}

%% ------------------------------------------------------------
\begin{frame}{Teste sua compreensão}

  \begin{quizbox}{}
    Um estimador por amostragem por importância é não-viesado para qualquer
    proposta $Q$ admissível. Então por que não coletamos um único lote gigante com
    $\pol_{\thv_0}$ e reutilizamos para sempre, reponderando?
    \paradiscussao
  \end{quizbox}

  \pause

  \begin{respbox}
    Porque \textbf{não-viesado não é o mesmo que útil}. À medida que $\thv$ se afasta
    de $\thv_0$, a variância cresce até que o tamanho amostral efetivo caia para
    ordem de uma ou duas trajetórias: o estimador continua correto \emph{em
    esperança}, mas cada realização é lixo. Os métodos de hoje atacam isso por dois
    lados --- \alert{limitar o quanto $\thv$ pode se afastar} (região de confiança,
    penalidade KL, recorte) e \alert{substituir a razão de trajetórias por uma razão
    de um único passo}.
  \end{respbox}
\end{frame}

%% ============================================================
\section{Objetivo substituto e PPO}
%% ============================================================

%% ------------------------------------------------------------
\begin{frame}{O objetivo substituto}
  \framesubtitle{Uma razão de um passo no lugar do produto}

  A identidade de desempenho relativo (Kakade e Langford, 2002):
  \[
    J(\polnovo) - J(\pol) = \frac{1}{1-\gamma}\,
      \E_{\substack{s\sim d^{\polnovo}\\ a\sim\polnovo}}\bigl[\Advpi(s,a)\bigr],
  \]
  com $d^{\pol}(s) = (1-\gamma)\sum_{t\ge0}\gamma^t \Prob(s_t = s\mid \pol)$, a \emph{ocupação descontada}.

  \vspace{-0.3em}
  \begin{alertabox}
    O incômodo: os estados vêm de $d^{\polnovo}$, política que ainda não executamos.
  \end{alertabox}

  \pause
  \vspace{-0.3em}
  \begin{defbox}{Objetivo substituto $\Lsub{\pol}(\polnovo)$}
    Trocamos $d^{\polnovo}$ por $\dpi$ (que \emph{temos}) e reponderamos apenas a ação:
    \[
      \Lsub{\pol}(\polnovo) \defeq \frac{1}{1-\gamma}
        \E_{\substack{s\sim \dpi\\ a\sim\pol}}
        \left[\frac{\polnovo(a\mid s)}{\pol(a\mid s)}\,\Advpi(s,a)\right].
    \]
    Uma razão \emph{por passo}, não um produto --- estimável com os dados de $\pol$.
  \end{defbox}
\end{frame}

%% ------------------------------------------------------------
\begin{frame}{PPO: duas maneiras de manter $\polnovo$ perto de $\pol$}

  \begin{teobox}{Variante 1 --- penalidade KL adaptativa}
    Resolve um problema \emph{irrestrito}:
    \[
      \thv_{k+1} = \argmax_{\thv}\; \Lsub{\thk}(\thv)
        - \beta_k\, \bar{D}_{\mathrm{KL}}(\thv \,\|\, \thk).
    \]
    O coeficiente $\beta_k$ é ajustado entre iterações: se a KL medida ficou acima do
    alvo, aumenta-se $\beta$; se ficou abaixo, diminui-se --- restrição KL aproximada.
  \end{teobox}

  \vspace{-0.4em}
  \begin{teobox}{Variante 2 --- objetivo recortado \emph{(a que usaremos)}}
    Seja $\razao = \dfrac{\polth(a_t\mid s_t)}{\pol_{\thk}(a_t\mid s_t)}$. Então
    \[
      \Lclip_{\thk}(\thv) = \E_{\traj\sim\pol_k}
        \Bigl[\textstyle\sum_{t=0}^{T}
          \min\bigl(\razao \Ahat_t,\;
          \clip(\razao,\,1-\epsilon,\,1+\epsilon)\,\Ahat_t\bigr)\Bigr],
    \]
    com $\epsilon$ típico $0{,}2$ e $\thv_{k+1} = \argmax_{\thv}\Lclip_{\thk}(\thv)$.
  \end{teobox}
\end{frame}

%% ------------------------------------------------------------
\begin{frame}{Teste sua compreensão}
  \framesubtitle{Os dois ramos do objetivo recortado}

  \[
    \Lclip = \min\bigl(r\,\Ahat,\; \clip(r,1-\epsilon,1+\epsilon)\,\Ahat\bigr),
    \qquad \epsilon\in(0,1)
  \]

  \begin{center}
  \scalebox{0.78}{\begin{tikzpicture}[font=\scriptsize]
    \foreach \dx/\rot in {0/A, 7.2/B}{
      \begin{scope}[xshift=\dx cm]
        \draw[rlCinza!60] (-0.15,-1.55) rectangle (4.6,1.55);
        \draw[->, rlCinza] (0,0) -- (4.4,0) node[below right] {$r$};
        \draw[->, rlCinza] (0,-1.45) -- (0,1.45);
        \node[font=\bfseries, rlAzul] at (2.2,1.85) {gráfico \rot};
      \end{scope}}
    % --- grafico A:  A>0  => L = min(r, 1+eps)  (subindo e depois plano)
    \begin{scope}
      \draw[rlAzul, very thick] (0,0) -- (2.9,1.16) -- (4.3,1.16);
      \draw[dashed, rlCinza] (1.9,-1.45) -- (1.9,0.76);
      \draw[dashed, rlCinza] (2.9,-1.45) -- (2.9,1.16);
      \node[below, rlCinza] at (1.9,-1.45) {$1-\epsilon$};
      \node[below, rlCinza] at (2.9,-1.45) {$1+\epsilon$};
      \node[below, rlCinza] at (2.4,-0.05) {$1$};
      \fill[rlLaranja] (2.4,0.96) circle (1.6pt);
    \end{scope}
    % --- grafico B:  A<0  => L = A*max(r,1-eps) (plano e depois descendo)
    \begin{scope}[xshift=7.2cm]
      \draw[rlAzul, very thick] (0,-0.76) -- (1.9,-0.76) -- (4.3,-1.72*0.8);
      \draw[dashed, rlCinza] (1.9,-1.45) -- (1.9,0.0);
      \draw[dashed, rlCinza] (2.9,-1.45) -- (2.9,0.0);
      \node[below, rlCinza] at (1.9,-1.45) {$1-\epsilon$};
      \node[below, rlCinza] at (2.9,-1.45) {$1+\epsilon$};
      \node[above, rlCinza] at (2.4,0.05) {$1$};
      \fill[rlLaranja] (2.4,-0.96) circle (1.6pt);
    \end{scope}
  \end{tikzpicture}}
  \end{center}

  \begin{quizbox}{}
    Qual dos dois gráficos corresponde a $\Adv > 0$ e qual a $\Adv < 0$?
    \paradiscussao
  \end{quizbox}
\end{frame}

%% ------------------------------------------------------------
\begin{frame}{Lendo os dois ramos}

  \begin{columns}[T]
  \begin{column}{0.49\textwidth}
    \begin{respbox}
      \textbf{Gráfico A: $\Adv>0$.} \; \textbf{Gráfico B: $\Adv<0$.}
    \end{respbox}
    \smallskip
    {\small
    \textbf{Se $\Adv>0$} (ação boa): queremos aumentar $r$. O objetivo cresce
    linearmente até $r = 1+\epsilon$ e \emph{satura}. Passado esse ponto o gradiente
    é zero: não há mais incentivo em empurrar a probabilidade da ação para cima.
    \par\smallskip
    \textbf{Se $\Adv<0$} (ação ruim): queremos diminuir $r$. O objetivo é constante
    para $r < 1-\epsilon$ --- já reduzimos o suficiente, pare. Mas do outro lado ele
    \emph{decresce sem limite}.}
  \end{column}
  \begin{column}{0.49\textwidth}
    \begin{alertabox}
      Duas leituras que costumam escapar:
      \begin{itemize}\setlength{\itemsep}{0.15ex}
        \item O recorte remove o \emph{incentivo} de ir mais longe; ele \textbf{não
          impede} que $r$ fique grande. Se a primeira época de otimização já jogou
          $r$ para $3$, o recorte não a traz de volta.
        \item Para $\Adv<0$ o objetivo é ilimitado inferiormente --- uma ação ruim
          pode ser punida arbitrariamente. É uma assimetria real do PPO.
      \end{itemize}
    \end{alertabox}
    \begin{ideiabox}
      O $\min$ faz de $\Lclip$ um limitante \emph{pessimista} (inferior) do
      substituto não recortado: nunca somos otimistas sobre o ganho de nos afastar.
    \end{ideiabox}
  \end{column}
  \end{columns}
\end{frame}

%% ------------------------------------------------------------
\begin{frame}{PPO recortado: o algoritmo}

  \begin{algbox}{PPO com objetivo recortado}
  \begin{itemize}\setlength{\itemsep}{0.1em}
    \item Inicialize os parâmetros da política $\thv_0$ e do crítico $w_0$.
    \item \textbf{para} $k = 0,1,2,\dots$ \textbf{faça}
      \begin{itemize}\setlength{\itemsep}{0.05em}
        \item Colete um lote $\Dtot_k$ de trajetórias executando $\pol_{\thk}$
          por $T$ passos, guardando $\log\pol_{\thk}(a_t\mid s_t)$.
        \item Calcule os retornos $\hat R_t$ e as vantagens $\Ahat_t$ (GAE
          truncado --- próxima seção); normalize $\Ahat$ no lote.
        \item \textbf{para} época $= 1,\dots,K$ (tipicamente $K \in [3,10]$) \textbf{faça}
          \begin{itemize}\setlength{\itemsep}{0.02em}
            \item Para cada minilote: suba o gradiente de
              $\Lclip_{\thk}(\thv) + c_{\ent}\,\ent[\polth] - c_v (\Vhat_w(s_t)-\hat R_t)^2$.
          \end{itemize}
        \item Interrompa cedo se $\bar{D}_{\mathrm{KL}}(\thv\|\thk)$ ultrapassar o alvo.
      \end{itemize}
  \end{itemize}
  \end{algbox}

  \begin{ideiabox}
    O ganho essencial está no laço interno: \textbf{várias atualizações de gradiente
    com o mesmo lote}. É de onde vem a eficiência amostral --- e é exatamente por
    isso que precisamos do recorte, porque a partir da segunda época os dados já são
    \emph{off-policy}.
  \end{ideiabox}
\end{frame}

%% ------------------------------------------------------------
\begin{frame}{PPO na prática}

  \begin{columns}[T]
  \begin{column}{0.5\textwidth}
    \textbf{Por que pegou}
    \begin{itemize}\setlength{\itemsep}{0.2ex}
      \item Eficiência de dados: $K$ passos de gradiente por lote coletado.
      \item Só primeiras derivadas --- sem inverter matriz de Fisher, sem gradiente
        conjugado (diferentemente do TRPO).
      \item Poucas linhas a mais que o REINFORCE.
      \item Robusto o bastante para virar padrão --- inclusive no ajuste de modelos
        de linguagem por RLHF.
    \end{itemize}
  \end{column}
  \begin{column}{0.5\textwidth}
    \textbf{O que ele não promete}
    \begin{itemize}\setlength{\itemsep}{0.2ex}
      \item Converge a ótimo \emph{local}, como todo gradiente de política.
      \item A melhoria monotônica é \emph{aproximada}: o recorte é uma heurística
        inspirada na teoria, não uma implementação dela.
      \item Sensível a detalhes de implementação (normalização de vantagem e de
        observação, recorte do gradiente, taxa decrescente, $\lambda$ do GAE).
    \end{itemize}
  \end{column}
  \end{columns}

  \begin{alertabox}
    A ideia de \emph{atualização conservadora de política} --- não confie no
    substituto longe do ponto onde ele foi construído --- é anterior ao PPO e
    remonta ao início dos anos 2000 (\emph{conservative policy iteration},
    Kakade e Langford). O PPO é a encarnação barata dessa ideia.
  \end{alertabox}
\end{frame}

%% ============================================================
\section{GAE: um botão para o viés e a variância}
%% ============================================================

%% ------------------------------------------------------------
\begin{frame}{Recapitulação: estimadores de $n$ passos e o erro TD}

  Da Aula 5, a família de estimadores de vantagem:
  \begin{align*}
    \Ahat^{(1)}_t &= r_t + \gamma V(s_{t+1}) - V(s_t)\\
    \Ahat^{(2)}_t &= r_t + \gamma r_{t+1} + \gamma^2 V(s_{t+2}) - V(s_t)\\
    \Ahat^{(\infty)}_t &= r_t + \gamma r_{t+1} + \gamma^2 r_{t+2} + \cdots - V(s_t)
  \end{align*}

  \begin{defbox}{Erro de diferença temporal}
    \[
      \delta^V_t \defeq r_t + \gamma V(s_{t+1}) - V(s_t)
    \]
    é o \emph{resíduo de Bellman} de um passo --- e, se $V = \Vpi$, um estimador
    não-viesado de $\Advpi(s_t,a_t)$.
  \end{defbox}

  \begin{ideiabox}
    Reparação central: $\Ahat^{(1)}_t = \delta^V_t$. Todo o resto desta seção é
    descobrir que os estimadores de $n$ passos são \emph{somas de $\delta$'s}.
  \end{ideiabox}
\end{frame}

%% ------------------------------------------------------------
\begin{frame}{Tudo é soma telescópica}

  \begin{align*}
    \Ahat^{(1)}_t &= \delta^V_t = r_t + \gamma V(s_{t+1}) - V(s_t)\\
    \Ahat^{(2)}_t &= \delta^V_t + \gamma\,\delta^V_{t+1}
      = r_t + \gamma r_{t+1} + \gamma^2 V(s_{t+2}) - V(s_t)\\
    \Ahat^{(k)}_t &= \sum_{l=0}^{k-1}\gamma^l \delta^V_{t+l}
      = \sum_{l=0}^{k-1}\gamma^l r_{t+l} + \gamma^k V(s_{t+k}) - V(s_t)
  \end{align*}

  \begin{center}
  \scalebox{0.85}{\begin{tikzpicture}[font=\scriptsize, node distance=3mm]
    \foreach \i in {0,...,4}{
      \node[draw=rlAzul, thick, fill=rlAzul!8, rounded corners=2pt,
            minimum width=15mm, minimum height=6mm] (d\i) at (\i*2.0,0)
            {$\gamma^{\i}\delta^V_{t+\i}$};}
    \node[rlCinza] at (5*2.0-0.2,0) {$\cdots$};
    \draw[trans destac] (-0.9,-0.75) -- (0.9,-0.75)
      node[midway, below, text=rlLaranja] {$\Ahat^{(1)}$};
    \draw[trans destac] (-0.9,-1.35) -- (2.9,-1.35)
      node[midway, below, text=rlLaranja] {$\Ahat^{(2)}$};
    \draw[trans destac] (-0.9,-1.95) -- (8.9,-1.95)
      node[midway, below, text=rlLaranja] {$\Ahat^{(\infty)}$: Monte Carlo};
  \end{tikzpicture}}
  \end{center}

  \vspace{-0.8em}
  \begin{ideiabox}
    Os $V$ internos se cancelam aos pares: sobram o \emph{bootstrapping} na ponta e a
    linha de base no início.
  \end{ideiabox}
\end{frame}

%% ------------------------------------------------------------
\begin{frame}{Estimador generalizado de vantagem (GAE)}

  Em vez de escolher \emph{um} $k$, tome a média exponencialmente ponderada de todos:
  {\small
  \begin{align*}
    \Ahat^{\mathrm{GAE}(\gamma,\lambda)}_t
      &= (1-\lambda)\bigl(\Ahat^{(1)}_t + \lambda \Ahat^{(2)}_t + \lambda^2\Ahat^{(3)}_t + \cdots\bigr)\\
      &= (1-\lambda)\bigl(\delta^V_t + \lambda(\delta^V_t + \gamma\delta^V_{t+1})
         + \lambda^2(\delta^V_t + \gamma\delta^V_{t+1} + \gamma^2\delta^V_{t+2}) + \cdots\bigr)\\
      &= (1-\lambda)\bigl(\delta^V_t(1+\lambda+\lambda^2+\cdots)
         + \gamma\delta^V_{t+1}(\lambda+\lambda^2+\cdots)
         + \gamma^2\delta^V_{t+2}(\lambda^2+\lambda^3+\cdots) + \cdots\bigr)\\
      &= (1-\lambda)\Bigl(\delta^V_t\tfrac{1}{1-\lambda}
         + \gamma\lambda\,\delta^V_{t+1}\tfrac{1}{1-\lambda}
         + \gamma^2\lambda^2\delta^V_{t+2}\tfrac{1}{1-\lambda} + \cdots\Bigr)\\[0.2em]
      &= \mdestaque{rlLaranja}{\sum_{l=0}^{\infty}(\gamma\lambda)^{l}\,\delta^V_{t+l}}
  \end{align*}}

  \vspace{-0.4em}
  \begin{ideiabox}
    Um desconto exponencial de razão $\gamma\lambda$ aplicado aos erros TD futuros.
    Schulman et al., \emph{High-Dimensional Continuous Control Using Generalized
    Advantage Estimation}, ICLR 2016.
  \end{ideiabox}
\end{frame}

%% ------------------------------------------------------------
\begin{frame}{Os dois extremos}

  \begin{columns}[T]
  \begin{column}{0.48\textwidth}
    \begin{teobox}{$\lambda = 0$}
      \[ \Ahat_t = \delta^V_t = r_t + \gamma V(s_{t+1}) - V(s_t) \]
      A vantagem de \textbf{TD(0)}. Usa uma recompensa real e o resto vem de $V$.
      \begin{itemize}\setlength{\itemsep}{0.1ex}
        \item Variância mínima.
        \item \alert{Viés} se $V \ne \Vpi$: o erro de $V(s_{t+1})$ entra no alvo.
      \end{itemize}
    \end{teobox}
  \end{column}
  \begin{column}{0.48\textwidth}
    \begin{teobox}{$\lambda = 1$}
      \[ \Ahat_t = \sum_{l\ge0}\gamma^l\delta^V_{t+l} = G_t - V(s_t) \]
      A vantagem de \textbf{Monte Carlo}. Os $V$ internos telescopam e somem.
      \begin{itemize}\setlength{\itemsep}{0.1ex}
        \item Variância máxima (soma de todas as recompensas aleatórias).
        \item $-V(s_t)$ é linha de base: não enviesa o gradiente nem se $V$ errar.
      \end{itemize}
    \end{teobox}
  \end{column}
  \end{columns}

  \begin{alertabox}
    Cuidado: $\lambda=1$ não dá estimador pontual não-viesado de $\Advpi$ ---
    $\E[\Ahat_t] = \Qpi(s_t,a_t) - V(s_t)$, igual a $\Advpi$ só se $V=\Vpi$. Não-viesado,
    para qualquer $V$, é o \textbf{gradiente}: a diferença é função só do estado.
  \end{alertabox}
\end{frame}

%% ------------------------------------------------------------
\begin{frame}{Teste sua compreensão}

  \[
    \Ahat^{\mathrm{GAE}(\gamma,\lambda)}_t = \sum_{l=0}^{\infty}(\gamma\lambda)^{l}\delta^V_{t+l},
    \qquad \delta^V_t = r_t + \gamma V(s_{t+1}) - V(s_t)
  \]

  \vspace{-0.9em}
  \begin{quizbox}{}
    Considere $\mathrm{GAE}(\gamma,0)$ e $\mathrm{GAE}(\gamma,0{,}99)$. Quais são verdadeiras?
    \begin{enumerate}\setlength{\itemsep}{0.1ex}\small
      \item[(a)] $\mathrm{GAE}(\gamma,0{,}99)$ é a vantagem calculada com um retorno TD(0).
      \item[(b)] $\mathrm{GAE}(\gamma,0)$ é a vantagem calculada com um retorno TD(0).
      \item[(c)] A \emph{variância} de $\mathrm{GAE}(\gamma,0)$ tende a ser maior.
      \item[(d)] O \emph{viés} de $\mathrm{GAE}(\gamma,0)$ tende a ser maior.
    \end{enumerate}
    \paradiscussao
  \end{quizbox}

  \pause
  \begin{respbox}
    Verdadeiras: \textbf{(b)} e \textbf{(d)}. A (a) troca os papéis; a (c) inverte a relação.
  \end{respbox}
\end{frame}

%% ------------------------------------------------------------
\begin{frame}{$\gamma$ e $\lambda$ fazem coisas diferentes}

  \begin{center}
  \scalebox{0.9}{\begin{tikzpicture}[font=\scriptsize]
    \draw[rlAzul!25, line width=8pt] (0,0) -- (10,0);
    \draw[trans destac, line width=1.6pt] (0,0) -- (10,0);
    \foreach \x/\lab in {0/{$\lambda=0$}, 5/{$\lambda \approx 0{,}95$}, 10/{$\lambda=1$}}{
      \fill[rlAzul] (\x,0) circle (2.4pt);
      \node[below=4mm, rlAzul, font=\scriptsize\bfseries] at (\x,0) {\lab};}
    \node[above=3mm, align=center, text=rlTeal] at (0,0)
      {TD(0)\\[-0.2em]{\tiny viés alto, variância baixa}};
    \node[above=3mm, align=center, text=rlLaranja] at (10,0)
      {Monte Carlo\\[-0.2em]{\tiny viés nulo, variância alta}};
    \node[above=3mm, align=center, text=rlVerde!65!black] at (5,0)
      {compromisso usual};
  \end{tikzpicture}}
  \end{center}

  \begin{columns}[T]
  \begin{column}{0.49\textwidth}
    \begin{defbox}{$\gamma$ --- desconto do \emph{problema}}
      Define qual objetivo estamos otimizando. Reduzi-lo introduz viés em relação
      ao problema original, mas o problema muda de fato.
    \end{defbox}
  \end{column}
  \begin{column}{0.49\textwidth}
    \begin{defbox}{$\lambda$ --- desconto do \emph{estimador}}
      Não muda o objetivo; muda apenas como estimamos a vantagem. Age sobre o quanto
      confiamos no crítico.
    \end{defbox}
  \end{column}
  \end{columns}

  \begin{alertabox}
    Os dois aparecem multiplicados ($\gamma\lambda$), o que os faz parecer
    intercambiáveis. Não são: $\gamma$ está no problema, $\lambda$ está no estimador.
  \end{alertabox}
\end{frame}

%% ------------------------------------------------------------
\begin{frame}{Medindo o compromisso}
  \framesubtitle{Corredor de 8 estados, $\gamma = 0{,}99$, política logística de um parâmetro}

  Gradiente exato $\partial J/\partial\thv = -0{,}2075$. Estimamos com
  $\mathrm{GAE}(\gamma,\lambda)$ sobre $30\,000$ trajetórias, variando o crítico:

  \begin{center}\small
  \begin{tabular}{@{}lrrrrrr@{}}
    \toprule
    & \multicolumn{2}{c}{$\Vhat = \Vpi$ (exato)}
    & \multicolumn{2}{c}{$\Vhat = 0{,}75\,\Vpi$}
    & \multicolumn{2}{c}{$\Vhat \equiv 0$}\\
    \cmidrule(lr){2-3}\cmidrule(lr){4-5}\cmidrule(lr){6-7}
    $\lambda$ & viés & desvio & viés & desvio & viés & desvio\\
    \midrule
    $0{,}00$ & $+0{,}002$ & $0{,}63$ & $+0{,}121$ & $0{,}84$ & $\mathbf{+0{,}478}$ & $2{,}25$\\
    $0{,}80$ & $+0{,}001$ & $0{,}90$ & $+0{,}097$ & $1{,}21$ & $+0{,}386$ & $3{,}29$\\
    $0{,}95$ & $+0{,}001$ & $1{,}12$ & $+0{,}033$ & $1{,}57$ & $+0{,}128$ & $4{,}41$\\
    $1{,}00$ & $+0{,}001$ & $1{,}23$ & $-0{,}005$ & $1{,}75$ & $-0{,}024$ & $5{,}03$\\
    \bottomrule
  \end{tabular}
  \end{center}

  \begin{alertabox}
    Olhe a célula em negrito. Sem crítico e com $\lambda=0$, o viés $+0{,}478$ é
    \emph{maior em módulo} que o próprio gradiente $-0{,}2075$: em média o estimador
    aponta para o \textbf{lado errado}. Nenhuma quantidade de amostras corrige isso.
  \end{alertabox}
\end{frame}

%% ------------------------------------------------------------
\begin{frame}{Por que $\lambda \approx 0{,}95$, e não $\lambda = 0$}

  Com um lote de $m$ trajetórias, o desvio cai como $1/\sqrt{m}$ --- mas o viés
  \textbf{não cai}. A raiz do erro quadrático médio,
  $\sqrt{\text{viés}^2 + \text{desvio}^2/m}$, com $\Vhat = 0{,}75\,\Vpi$:

  \begin{center}\small
  \begin{tabular}{@{}lrrrr@{}}
    \toprule
    $\lambda$ & $m=1$ & $m=10$ & $m=100$ & $m=1000$\\
    \midrule
    $0{,}00$ & $\mathbf{0{,}849}$ & $\mathbf{0{,}292}$ & $\mathbf{0{,}147}$ & $0{,}124$\\
    $0{,}80$ & $1{,}219$ & $0{,}396$ & $0{,}155$ & $0{,}104$\\
    $0{,}95$ & $1{,}566$ & $0{,}496$ & $0{,}160$ & $0{,}059$\\
    $0{,}98$ & $1{,}672$ & $0{,}529$ & $0{,}168$ & $\mathbf{0{,}054}$\\
    \bottomrule
  \end{tabular}
  \end{center}

  \begin{ideiabox}
    \textbf{O $\lambda$ ótimo cresce com o tamanho do lote.} Com poucas trajetórias a
    variância domina e vale a pena aceitar viés ($\lambda$ pequeno); com lotes
    grandes a variância já está controlada e o viés passa a ser o gargalo. Os
    valores usuais na literatura ($\lambda \in [0{,}9;\,0{,}99]$) são a resposta certa
    para o regime em que o PPO opera: milhares de passos por atualização.
  \end{ideiabox}
\end{frame}

%% ------------------------------------------------------------
\begin{frame}{GAE truncado: como o PPO realmente calcula}

  A soma infinita é impossível; corta-se no fim do lote de $T$ passos:
  \[
    \Ahat_t = \sum_{l=0}^{T-t-1}(\gamma\lambda)^l\,\delta^V_{t+l}.
  \]

  \begin{algbox}{Cálculo em uma varredura reversa (custo $\mathcal{O}(T)$)}
  \begin{itemize}\setlength{\itemsep}{0.1em}
    \item $\Ahat_{T} \leftarrow 0$
    \item \textbf{para} $t = T-1$ até $0$ \textbf{faça}
      \begin{itemize}
        \item $\delta_t \leftarrow r_t + \gamma\,\Vhat(s_{t+1})\,(1-\mathrm{fim}_t) - \Vhat(s_t)$
        \item $\Ahat_t \leftarrow \delta_t + \gamma\lambda\,(1-\mathrm{fim}_t)\,\Ahat_{t+1}$
      \end{itemize}
    \item Alvo do crítico: $\hat R_t \leftarrow \Ahat_t + \Vhat(s_t)$.
  \end{itemize}
  \end{algbox}

  \begin{itemize}\setlength{\itemsep}{0.2ex}
    \item \textbf{Benefício:} atualiza-se a cada $T$ passos, sem esperar o fim do episódio.
    \item O truncamento introduz viés (a cauda além de $T$ some), pequeno se
      $(\gamma\lambda)^{T}\ll1$. O indicador $\mathrm{fim}_t$ zera o \emph{bootstrapping}
      em estados terminais --- erro clássico de implementação.
  \end{itemize}
\end{frame}

%% ============================================================
\section{Por que isso funciona: melhoria monotônica}
%% ============================================================

%% ------------------------------------------------------------
\begin{frame}{Voltando ao limite entre duas políticas}

  Da aula passada: usamos $\dpi$ no lugar de $d^{\polnovo}$. \emph{Por quê?}
  Porque não temos amostras de $d^{\polnovo}$. Mas a troca custa alguma coisa ---
  e a teoria diz exatamente quanto.

  \begin{teobox}{Limite de desempenho relativo (Achiam, Held, Tamar, Abbeel, 2017)}
    \[
      \Bigl\lvert\, J(\polnovo) - \bigl(J(\pol) + \Lsub{\pol}(\polnovo)\bigr) \Bigr\rvert
      \;\le\; C \sqrt{\;\E_{s\sim \dpi}\bigl[\Dkl(\polnovo \,\|\, \pol)[s]\bigr]\;},
    \]
    com $\Lsub{\pol}(\polnovo) = \frac{1}{1-\gamma}
      \E_{s\sim\dpi,\,a\sim\pol}\bigl[\frac{\polnovo(a\mid s)}{\pol(a\mid s)}\Advpi(s,a)\bigr]$
    e $C$ proporcional a $\dfrac{\gamma \max_s\lvert\E_{a\sim\polnovo}\Advpi(s,a)\rvert}{(1-\gamma)^2}$.
  \end{teobox}

  \begin{ideiabox}
    Leitura em uma frase: \emph{o substituto é confiável exatamente na medida em que
    as duas políticas estão próximas em KL}. A aproximação não é ruim --- ela é
    ruim \emph{longe}, com erro controlado.
  \end{ideiabox}
\end{frame}

%% ------------------------------------------------------------
\begin{frame}{De um limite a um algoritmo}

  Tirando o valor absoluto e ficando com o lado que interessa:
  \[
    J(\polnovo) - J(\pol) \;\ge\;
      \mdestaque{rlVerde}{\Lsub{\pol}(\polnovo) - C\sqrt{\E_{s\sim\dpi}[\Dkl(\polnovo\|\pol)[s]]}} .
  \]

  \begin{columns}[T]
  \begin{column}{0.52\textwidth}
    \begin{itemize}\setlength{\itemsep}{0.3ex}
      \item Se \textbf{maximizarmos o lado direito} em $\polnovo$, garantimos
        melhoria do lado esquerdo.
      \item É um algoritmo de \alert{majorizar--maximizar} (MM) aplicado ao
        objetivo verdadeiro.
      \item E o milagre prático: \emph{os dois} termos do lado direito são
        estimáveis com amostras de $\pol$ --- a política que já rodamos.
    \end{itemize}
  \end{column}
  \begin{column}{0.46\textwidth}
    \begin{center}
    \scalebox{0.82}{\begin{tikzpicture}[font=\scriptsize]
      \draw[->, rlCinza] (-0.2,0) -- (5.4,0) node[below] {$\thv$};
      \draw[->, rlCinza] (-0.2,0) -- (-0.2,3.3) node[left] {$J$};
      % objetivo verdadeiro
      \draw[rlAzul, very thick, domain=0:5.2, samples=90]
        plot (\x, {1.1 + 1.6*sin(\x*38) *0.55 + 0.30*\x});
      \node[rlAzul, font=\scriptsize\bfseries] at (4.9,3.05) {$J(\thv)$};
      % minorante em torno de theta_k = 1.5
      \draw[rlVerde!70!black, very thick, dashed, domain=0.2:3.0, samples=70]
        plot (\x, {1.1 + 1.6*sin(1.5*38)*0.55 + 0.30*1.5 - 0.75*(\x-1.5)^2 + 0.62*(\x-1.5)});
      \node[rlVerde!60!black, font=\scriptsize\bfseries, align=center] at (1.1,0.42)
        {minorante\\ substituto};
      \fill[rlLaranja] (1.5,{1.1 + 1.6*sin(1.5*38)*0.55 + 0.45}) circle (2.2pt);
      \node[rlLaranja, font=\scriptsize] at (1.5,{1.1 + 1.6*sin(1.5*38)*0.55 + 0.95}) {$\thk$};
      \fill[rlVermelho] (1.91,{1.1 + 1.6*sin(1.91*38)*0.55 + 0.573}) circle (2.2pt);
      \node[rlVermelho, font=\scriptsize] at (2.55,{1.1 + 1.6*sin(1.91*38)*0.55 + 0.30}) {$\thv_{k+1}$};
    \end{tikzpicture}}
    \end{center}
  \end{column}
  \end{columns}

  \begin{ideiabox}
    Mesmo esqueleto de EM e do ponto proximal: minorante que toca a função no ponto
    atual, maximize-o, repita.
  \end{ideiabox}
\end{frame}

%% ------------------------------------------------------------
\begin{frame}{A garantia de melhoria}

  \begin{teobox}{Proposição}
    Se $\polk$ e $\pol_{k+1}$ estão relacionadas por
    \[
      \pol_{k+1} = \argmax_{\polnovo}\;
        \Lsub{\polk}(\polnovo) - C\sqrt{\E_{s\sim d^{\polk}}[\Dkl(\polnovo\|\polk)[s]]},
    \]
    então $J(\pol_{k+1}) \ge J(\polk)$.
  \end{teobox}

  \textbf{Demonstração.} Basta exibir um ponto viável com valor $0$:
  \begin{enumerate}\setlength{\itemsep}{0ex}
    \item $\polk$ é viável, e $\Dkl(\polk\|\polk)[s] = 0$ para todo $s$.
    \item $\Lsub{\polk}(\polk) = \frac{1}{1-\gamma}\E_{s\sim d^{\polk},a\sim\polk}[A^{\polk}(s,a)] = 0$,
      pois a vantagem tem média zero sob a própria política.
    \item Logo o objetivo em $\polk$ vale $0$ e o \emph{ótimo} é $\ge 0$; pelo limite
      anterior, $J(\pol_{k+1}) - J(\polk) \ge 0$. \hfill$\blacksquare$
  \end{enumerate}

  \begin{ideiabox}
    A demonstração sobrevive se a otimização for restrita a uma classe parametrizada
    $\Pi_{\thv}$, desde que $\polk \in \Pi_{\thv}$ --- por isso vale para redes neurais.
  \end{ideiabox}
\end{frame}

%% ------------------------------------------------------------
\begin{frame}{Por que ninguém usa esse algoritmo literalmente}

  \[
    \pol_{k+1} = \argmax_{\polnovo}\;
      \Lsub{\polk}(\polnovo) - C\sqrt{\E_{s\sim d^{\polk}}[\Dkl(\polnovo\|\polk)[s]]}
  \]

  \begin{alertabox}
    O $C$ que a teoria fornece cresce como $\gamma/(1-\gamma)^2$. Com
    $\gamma = 0{,}99$ isso é da ordem de $10^4$: a penalidade domina o substituto e
    os passos ficam \emph{absurdamente pequenos}. A garantia é verdadeira e inútil.
  \end{alertabox}

  \begin{ideiabox}
    Duas saídas, e as duas abandonam a constante teórica: transformar a penalidade em
    \textbf{restrição} com um raio $\delta$ escolhido à mão (região de confiança,
    TRPO), ou \textbf{ajustar} o coeficiente da penalidade por realimentação (PPO).
  \end{ideiabox}
\end{frame}

%% ------------------------------------------------------------
\begin{frame}{Três aproximações do mesmo algoritmo idealizado}

  \begin{center}\small
  \begin{tabular}{@{}p{0.30\linewidth}p{0.30\linewidth}p{0.32\linewidth}@{}}
    \toprule
    \textbf{TRPO} & \textbf{PPO--penalidade} & \textbf{PPO--recorte}\\
    \midrule
    Substitui a penalidade por uma \emph{restrição} rígida
      $\bar{D}_{\mathrm{KL}} \le \delta$ &
    Mantém a penalidade, mas \emph{ajusta} $\beta_k$ para atingir uma KL alvo &
    Abandona a KL: recorta a razão em $[1-\epsilon, 1+\epsilon]$\\[0.4em]
    Passo natural + busca em linha; precisa de gradiente conjugado &
    Só primeiras derivadas; um hiperparâmetro extra &
    Só primeiras derivadas; um hiperparâmetro ($\epsilon$)\\[0.4em]
    Caro por iteração &
    Intermediário &
    Barato e simples de implementar\\
    \bottomrule
  \end{tabular}
  \end{center}

  \begin{ideiabox}
    A teoria da melhoria monotônica não é \emph{executada} por nenhum dos três. Ela
    explica \emph{por que} vale a pena restringir o passo --- e é essa explicação, não
    a fórmula, que sobreviveu à prática.
  \end{ideiabox}
\end{frame}

%% ------------------------------------------------------------
\begin{frame}{Teste sua compreensão}

  \begin{quizbox}{}
    Sobre o REINFORCE, quais afirmações são verdadeiras?
    \begin{enumerate}\setlength{\itemsep}{0.1ex}\small
      \item[(a)] Acrescentar uma linha de base ajuda a reduzir a variância das
        atualizações.
      \item[(b)] Ele converge para um ótimo global.
      \item[(c)] Pode ser inicializado com uma política determinística subótima e
        ainda assim convergir para um ótimo local, dados passos adequados.
      \item[(d)] Um único passo de gradiente pode produzir uma política
        \emph{pior} que a inicial em retorno obtido.
    \end{enumerate}
    \paradiscussao
  \end{quizbox}

  \pause
  \begin{respbox}
    \textbf{(a)} e \textbf{(d)} são verdadeiras. \textbf{(b)} é falsa: a otimização é
    não-convexa e converge a ótimos locais. \textbf{(c)} é falsa e vale insistir: uma
    política \emph{determinística} nunca toma as demais ações --- sem exploração, o
    gradiente não enxerga as alternativas. \textbf{(d)} é a motivação desta seção.
  \end{respbox}
\end{frame}

%% ------------------------------------------------------------
\begin{frame}{Balanço da aula}

  \begin{columns}[T]
  \begin{column}{0.49\textwidth}
    \begin{defbox}{Gradientes de política --- o que ficou}
      \begin{itemize}\setlength{\itemsep}{0.1ex}\small
        \item Família muito usada, bem além deste curso.
        \item Funciona quando a recompensa \textbf{não é diferenciável}.
        \item Quase sempre combinada com valor sem modelo: \emph{ator--crítico}.
        \item O crítico entra por duas portas --- linha de base (variância) e
          \emph{bootstrapping} (viés); o GAE interpola entre elas.
      \end{itemize}
    \end{defbox}
  \end{column}
  \begin{column}{0.49\textwidth}
    \begin{center}
    \scalebox{0.78}{\begin{tikzpicture}[font=\scriptsize, node distance=4mm]
      \node[draw=rlCinza, thick, fill=rlFundo, rounded corners=2pt,
            minimum width=30mm, minimum height=7mm] (a) {REINFORCE};
      \node[draw=rlAzul, thick, fill=rlAzul!8, rounded corners=2pt,
            minimum width=30mm, minimum height=7mm, below=5mm of a] (b) {$+$ linha de base};
      \node[draw=rlTeal, thick, fill=rlTeal!10, rounded corners=2pt,
            minimum width=30mm, minimum height=7mm, below=5mm of b] (c) {ator--crítico $+$ GAE};
      \node[draw=rlAmbar, thick, fill=rlAmbar!12, rounded corners=2pt,
            minimum width=30mm, minimum height=7mm, below=5mm of c] (d) {TRPO};
      \node[draw=rlLaranja, very thick, fill=rlLaranja!14, rounded corners=2pt,
            minimum width=30mm, minimum height=7mm, below=5mm of d] (e) {PPO};
      \draw[trans] (a)--(b); \draw[trans] (b)--(c);
      \draw[trans] (c)--(d); \draw[trans destac] (d)--(e);
      \node[right=3mm of b, text=rlCinza, font=\tiny, align=left] {menos variância};
      \node[right=3mm of c, text=rlCinza, font=\tiny, align=left] {viés controlado};
      \node[right=3mm of d, text=rlCinza, font=\tiny, align=left] {passo controlado};
      \node[right=3mm of e, text=rlCinza, font=\tiny, align=left] {e barato};
    \end{tikzpicture}}
    \end{center}
  \end{column}
  \end{columns}

  \begin{alertabox}
    Tudo isso pressupõe que existe uma \textbf{função de recompensa}. E se não
    existir --- se tudo o que temos for alguém fazendo a tarefa corretamente?
    É o ponto de partida da Aula 8.
  \end{alertabox}
\end{frame}

%% ============================================================
\section{Síntese}
%% ============================================================

%% ------------------------------------------------------------
\begin{frame}{O que você deve levar desta aula}

  \begin{defbox}{Capacidades esperadas}
    \begin{itemize}\setlength{\itemsep}{0.2ex}
      \item Escrever a identidade de amostragem por importância e explicar por que a
        \emph{variância}, não o viés, é o obstáculo.
      \item Mostrar que a razão de trajetórias colapsa em um produto de razões de
        política --- e por que esse produto é instável.
      \item Enunciar as duas variantes do PPO e \textbf{ler o gráfico do recorte}
        nos dois sinais da vantagem.
      \item Derivar $\Ahat^{\mathrm{GAE}(\gamma,\lambda)} = \sum_l (\gamma\lambda)^l\delta^V_{t+l}$
        e identificar os casos $\lambda = 0$ e $\lambda = 1$.
      \item Demonstrar a garantia de melhoria monotônica a partir do limite de
        desempenho relativo, e dizer por que o $C$ da teoria inviabiliza o uso literal.
    \end{itemize}
  \end{defbox}
\end{frame}

%% ------------------------------------------------------------
\begin{frame}{Leituras e próxima aula}

  \begin{columns}[T]
  \begin{column}{0.5\textwidth}
    \textbf{Leituras}
    \begin{itemize}\setlength{\itemsep}{0.2ex}\small
      \item Schulman et al., \emph{Generalized Advantage Estimation}, ICLR 2016.
      \item Schulman et al., \emph{Proximal Policy Optimization Algorithms}, 2017.
      \item Achiam, Held, Tamar e Abbeel, \emph{Constrained Policy Optimization}, ICML 2017.
    \end{itemize}
  \end{column}
  \begin{column}{0.48\textwidth}
    \begin{ideiabox}
      \textbf{Próxima aula.} Aprender a partir de demonstrações e de preferências:
      clonagem de comportamento, DAgger, RL inverso, RLHF e DPO.
    \end{ideiabox}
    \begin{alertabox}
      Uma pergunta para levar: por que o recorte do PPO funciona tão bem na
      prática, se não garante melhoria monotônica?
    \end{alertabox}
  \end{column}
  \end{columns}
\end{frame}

\end{document}
